\section{Related Work}
\label{sec:relatedwork}

We organise related work around the novelty threats that matter for this paper
rather than as a catalogue.

\paragraph{Quantisation and JHQ.}
Product Quantization (PQ)~\cite{pq} established Cartesian decomposition across
subspaces and asymmetric distance computation. JHQ~\cite{jhq} targets
high-dimensional ANN with an orthogonal transform, a Cartesian primary
quantiser, and hierarchical residual refinement. We do not claim JHQ's
hierarchy or orthogonal transform as contributions; our work begins from that
representation and asks how its specific structure should execute on a GPU.

\paragraph{Grouped lookup tables.}
This is the closest prior work and we state the overlap plainly. PQ Fast
Scan~\cite{fastscan} showed that cache residency alone is insufficient for fast
ADC; Quicker ADC~\cite{quickeradc} generalised it with irregular product
quantisers and split tables. Most directly, GPU IVF-RaBitQ~\cite{ivfrabitq}
partitions dimensions into groups of $U$ bits, builds one $2^{U}$-entry table
per group, sums $D/U$ lookups per candidate, and selects $U{=}4$ by trading
shared-memory capacity against arithmetic. That trade-off is the same one
Section~\ref{sec:designspace} navigates, it is prior art, and we do not claim
it.

Two things are ours. First, the object: RaBitQ groups a representation designed
to be grouped --- its code \emph{is} the bit pattern its table is indexed by ---
whereas JHQ's primary code is an 8-bit index into 256 codewords with no grouped
structure on its face. Equation~\ref{eq:split} is an identity that has to be
derived from the codebook's Cartesian construction, and it does not exist for an
ordinary $k$-means product quantiser at the same code width. Second, and more
useful, the method: RaBitQ selects its group size by measuring which performs
best, while Section~\ref{sec:rq3} decomposes the choice into an issue term read
off the disassembled kernel and a residency term measured as the residual
against it. That decomposition also produces a result the tuned-choice approach
would not surface --- at JHQ's code width the 16-entry tables that RaBitQ
settles on ($G{=}4$ here) are $4$--$36\%$ slower than 32-entry ones, so the
optimal group size is a property of the code width and not a constant.

\paragraph{Adaptive ANN control.}
ANN libraries such as FAISS~\cite{faiss} commonly expose parameter sweeps and
tuning utilities, while learned approaches can adapt search termination to the
query~\cite{earlyterm}. We therefore do not claim sampling or parameter
selection as new. Our output-stability rule differs in its signal and deployment
assumptions: it needs no labelled nearest neighbours and no learned predictor,
compares the system's own returned top-$k$ sets to a generous-budget reference,
and chooses the externally specified JHQ refinement budget. The evaluation
separately quantifies selection error and amortisation, which are necessary
because a label-free proxy is not a recall guarantee.

\paragraph{GPU ANN systems.}
CAGRA~\cite{cagra} represents the graph-based GPU frontier, and cuVS provides
production implementations of CAGRA, IVF-PQ, and IVF-RaBitQ. GPU
IVF-RaBitQ~\cite{ivfrabitq} combines inverted-file routing with low-bit RaBitQ
codes and reports a strong recall--throughput/build/storage trade-off. These
systems are not uniformly weaker than JHQ: our matched experiments show regime
crossings in the high-recall tail and with batch size. This is why we report
nondominated frontiers and configuration-specific build failures rather than a
single aggregate speedup.
